\chapter{Introduction}

Since Knuth popularized the term \textit{analysis of algorithms} \cite[Preface]{knuth1997art}, the field has grown to include highly specialized subfields. One useful way to classify such analyses is by two aspects. The first is the scenario under which the performance of an algorithm is analyzed, such as worst-case or average-case analysis. \cyril{``in general'' is too strong, change it into something like ``in many cases''.}\minh{I agree} In many cases, average-case analysis is more informative than worst-case analysis, and is often more challenging to perform, since it requires a probabilistic model of the input that is both tractable and expressive enough. For example, in the case of regular languages, generating syntax trees uniformly at random may be unsuitable \cite{koechlin2019uniform}. Average-case analysis also opens the door to techniques from other fields, such as probability theory and analytic combinatorics, and raises new questions, including questions about limiting objects as the input size tends to infinity. The second aspect is the underlying cost model, such as the uniform model or the word-RAM model. A cost model should align quantitative analysis with practical performance, for example execution time. The analysis of modern algorithms may therefore require cost models that capture hardware features such as caching, branch prediction, and parallelism.

In this internship, we study Hyperscan, a regular expression matching engine that uses SIMD parallelism to achieve high performance. Because Hyperscan has a sophisticated implementation, we focus on questions that can be isolated and studied mathematically. Hyperscan extracts possible strings from a regular expression and matches them in parallel using SIMD string matchers. It is therefore natural to analyze statistics related to those strings under a distribution. In particular, our main question is to understand the asymptotics of the length of the longest string extracted from a regular expression. We also formulate related questions that were not studied in depth. Details are given in the next chapter.

The rest of the report is organized as follows. \textbf{Chapter 2} introduces the concepts and techniques used in this internship, including regular expressions, BST-like distributions, SIMD parallelism, analytic combinatorics, and Monte Carlo simulation. In \textbf{Chapter 3}, we present Hyperscan and several problems arising from its analysis. \textbf{Chapter 4} contains our results on the length of the longest extracted string. Where useful, proofs are preceded by the motivation behind the approach. Finally, \textbf{Chapter 5} summarizes the work and discusses future directions.
